\chapter{Conclusion and Future Work}
\label{ch:conclusion}

This work set out to determine whether a small, publicly available language model, taught the language of mathematics by supervised fine-tuning, can be steered by reinforcement learning to discover symbolic expressions --- and, more pointedly, \emph{which parts of the RL recipe actually earn their cost}. The Seriguela pipeline provided the laboratory: a 682K-expression synthetic corpus, six fine-tuned GPT-2 variants, and a discovery loop in which five search strategies could be compared under identical rewards, budgets, and seeds. This chapter revisits the research questions in light of the results, discusses what the findings mean beyond the benchmarks, acknowledges the limitations of the study, and outlines the paths it opens.

\section{Answers to the Research Questions}

\textbf{RQ1 --- Does RL beat sampling at the same budget?} Yes, but only where it matters. On saturated problems the fine-tuned prior alone reaches $R^2 \approx 0.99$ and learning adds nothing. On the needle-in-a-haystack problem Nguyen-5, the best RL arm adds $+0.23$ mean best $R^2$ and $+80$ percentage points of exact-discovery rate over Best-of-N at an identical sampling budget. RL for symbolic regression is not a blanket improvement; it is an investment that pays precisely on the problems the prior cannot already solve.

\textbf{RQ2 --- Does the elite buffer help?} Conditionally --- and the condition is the finding. The buffer transformed PPO (from 2/5 to 4/5 exact solutions on Nguyen-5) and did nothing for GRPO (BoN-GRPO $\approx$ Pure-GRPO everywhere). Retention only creates value when exploration first produces something worth retaining. This conditional effect, predicted by the exploration--retention decomposition and confirmed across all four problems, is the thesis's central empirical contribution.

\textbf{RQ3 --- PPO or GRPO?} They differ less in final scores than in the \emph{search regimes} their updates induce. PPO, once tuned below its collapse threshold, operates in a meta-stable regime near the edge of syntactic instability, exploring roughly ten times more unique expressions than GRPO; GRPO is remarkably stable at default hyperparameters --- it was the only family that required no tuning at all --- but its group-normalized advantages lose selective pressure when rewards are homogeneous, freezing the search within about fifteen steps. For symbolic discovery, where the goal is a rare exact hit rather than a good average, PPO's regime wins; for applications valuing stability and syntactic cleanliness over peak discovery, GRPO's profile may be preferable.

\textbf{RQ4 --- Does the model generalize beyond its dimensional training scope?} Approximatively, yes. Median performance declines smoothly with dimensionality (0.997 at one variable down to 0.623 at five) and does not collapse at the training boundary: on problems with six to nine variables --- more than any expression the model was fine-tuned on --- it still reaches best $R^2$ between 0.22 and 0.79 (median 0.454 beyond scope versus 0.863 within). No exact solution was found beyond five variables, so the claim is graceful degradation of \emph{approximation}, not extrapolation of \emph{symbolic discovery}.

\textbf{RQ5 --- How does it compare to the state of the art?} Under the SRBench ground-truth protocol, the champion solves 16.2\% of the Feynman problems exactly --- essentially matching gplearn and within half a point of Operon on the published leaderboard, with three structural handicaps held against it: frozen constants ($C=1$), a 124M-parameter model, and a fixed 200-step budget. The result should be read not as a leaderboard claim but as an existence proof: a small language model with a reused pre-trained prior is a viable symbolic regression engine, and the constant-fitting ablation indicates the reported numbers are a \fic{lower bound, with roughly three points of solution rate recoverable by post-hoc fitting}. % FICTITIOUS: ablation numbers provisional

\section{Interpretability and Scientific Utility}

Symbolic regression's promise is a model a scientist can read. The expressions the champion produced keep that promise in two ways. First, they are short: across 122 problems, discovered expressions cluster at the complexity scale of the targets themselves, with no bloat --- notably, without any parsimony term in the reward, an effect we attribute to the short-expression prior inherited from the SFT corpus. Second, they are structurally meaningful: the model rewrites $x^y$ as $e^{y \ln x}$, composes $\cos/\tan$ to express a cotangent missing from its vocabulary, and reaches the same trigonometric identity through several algebraically equivalent syntactic routes. These are the behaviors of a system that has internalized compositional structure, not memorized a lookup table --- the property that would let such a system propose \emph{candidate laws} over experimental data in a real discovery setting.

\section{Limitations}

The study's boundaries should be stated as plainly as its findings.

\begin{itemize}
    \item \textbf{Frozen constants.} The $C=1$ protocol, chosen to keep the reward uniform and cheap, handicaps every accuracy metric relative to methods that fit constants. The ablation of Section~\ref{sec:constant_ablation} \fic{quantifies the cost at a few points of solution rate}, but constant handling \emph{inside} the RL loop remains untested. % FICTITIOUS: ablation provisional
    \item \textbf{One model, one notation.} The RL phases ran exclusively on GPT-2 Base Infix. The zero-shot screening justified the scale choice empirically, but the notation choice (infix over prefix) rests on preliminary evidence only, and conclusions may not transfer to other architectures.
    \item \textbf{Narrow discriminative base.} The algorithm ranking rests heavily on Nguyen-5 --- the one problem in the Phase 2 set hard enough to separate the arms. The Phase 4 sweep corroborates the champion at scale, but does not re-run the five-arm comparison there; a comparison on more needle-class problems would strengthen the mechanism's generality.
    \item \textbf{Statistical power.} Five seeds per cell detect only large effects ($d \gtrsim 1.5$); the PPO arms' bimodal outcomes further violate ANOVA assumptions, which is why discovery rates carry the primary evidential weight. Small-to-moderate differences between arms remain open. Phase 4, moreover, ran a single seed per problem: its per-problem scores are point estimates, and its aggregate rates carry sampling noise that replication would reduce.
    \item \textbf{Asymmetric tuning.} PPO required hyperparameter tuning to avoid collapse; GRPO ran on defaults. The comparison is conservative in GRPO's favor, but a fully symmetric tuning budget for both families was beyond the compute available.
    \item \textbf{Solution criterion.} Our solution rate uses test $R^2 > 0.999$ as a proxy for SRBench's SymPy-verified symbolic equivalence; spot checks confirmed agreement on the champion's solutions, but the automated criterion can in principle over- or under-count edge cases.
    \item \textbf{Clean data only.} Noise injection and adversarial prompt conditioning are implemented in the framework but were fixed (noise 0, standard prompt) in the final protocol; robustness claims are therefore out of scope.
\end{itemize}

\section{Future Work}

Several directions follow naturally from the results.

\textbf{Constants in the loop.} The single largest accuracy lever is constant handling. A curriculum that begins with $C=1$ and introduces cheap constant fitting (e.g., a few L-BFGS-B iterations on the batch elite) once validity stabilizes would test whether the exploration--retention dynamics survive a denser, costlier reward.

\textbf{Modern small models.} The pipeline is model-agnostic: repeating the comparison with contemporary small language models (e.g., Llama- or Qwen-family models in the 0.5--3B range) would test whether stronger pre-trained priors shift the balance between sampling and learning --- and whether the PPO collapse threshold moves with model quality.

\textbf{Exploiting the implemented robustness axes.} The noise-injection and prompt-conditioning machinery (oracle and distractor operator lists) is built and documented; a robustness study over these factors is the cheapest substantial extension of this work, and directly relevant to real experimental data.

\textbf{Retention beyond the buffer.} The elite buffer is the simplest member of a family that includes priority-queue training, replay with learned priorities, and population-based methods. Given that retention was the binding constraint for PPO, richer retention mechanisms --- e.g., buffers stratified by expression skeleton --- may raise the exact-discovery ceiling further.

\textbf{Toward scientific practice.} The end goal is deployment on real measurement data: physics and chemistry datasets with units, noise, and missing variables, where the human-in-the-loop interaction pioneered by systems like Sym-Q \cite{tian_interactive_2025} could combine the model's structural prior with an expert's domain constraints.

\section{Final Remarks}

The contribution of this thesis is ultimately a decomposition. ``Reinforcement learning for symbolic regression'' is not one ingredient but two --- a search regime that must be wide enough to find rare structure, and a memory that must be reliable enough to keep it --- and the popular algorithms distribute these ingredients unevenly. Naming the factors, measuring them separately, and showing that only their conjunction solves the hard problems turns a leaderboard observation into a mechanism. We hope the pipeline, the released artifacts, and the failure catalogue documented along the way lower the cost for the next investigation of this kind.
